\section*{NeurIPS Paper Checklist}

\begin{enumerate}

\item {\bf Claims}
    \item[] Question: Do the main claims made in the abstract and introduction accurately reflect the paper's contributions and scope?
    \item[] Answer: \answerYes{}
    \item[] Justification: The abstract and introduction match the proved balanced $R=2$ theorem, the dustbin bridge, the local surrogate-bias bound, the exactness checks, and the reported TPU runs; they do not claim a full KL-unbalanced theory or CE-optimized promoted training.

\item {\bf Limitations}
    \item[] Question: Does the paper discuss the limitations of the work performed by the authors?
    \item[] Answer: \answerYes{}
    \item[] Justification: The discussion explicitly limits the theorem to fixed-support balanced transport, states that the dustbin bridge is not a full unbalanced OT theorem, and notes that the promoted run optimizes reconstruction while CE is diagnostic only.

\item {\bf Theory assumptions and proofs}
    \item[] Question: For each theoretical result, does the paper provide the full set of assumptions and a complete (and correct) proof?
    \item[] Answer: \answerYes{}
    \item[] Justification: The assumptions and statements for the balanced $R=2$ adjoint, the one-reference-tile backward schedule, the dustbin augmentation bridge, the local surrogate-bias bound, the projective contraction certificate, the a posteriori bias certificate, the certified tail-depth optimization rule, and the transport-orbit calculus are given in Section 2, with complete proofs in the appendices.

\item {\bf Experimental result reproducibility}
    \item[] Question: Does the paper fully disclose all the information needed to reproduce the main experimental results of the paper to the extent that it affects the main claims and/or conclusions of the paper (regardless of whether the code and data are provided or not)?
    \item[] Answer: \answerYes{}
    \item[] Justification: The manuscript discloses the exact surrogate, kernel schedule, validation settings, TPU type, dataset family, screen configurations, promoted-run objective weights, throughput, checkpoint steps, and held-out evaluation protocol. The anonymized supplementary ZIP contains the analyzed transport kernel and a Sinkhorn-only Pfam training harness with the reported v6e-8/default settings; it omits cloud launchers, project-specific storage paths, checkpoints, logs, and external data or embedding generation.

\item {\bf Open access to data and code}
    \item[] Question: Does the paper provide open access to the data and code, with sufficient instructions to faithfully reproduce the main experimental results, as described in supplemental material?
    \item[] Answer: \answerYes{}
    \item[] Justification: An anonymized supplementary bundle accompanies the submission and contains the single-file kernel implementation for the balanced and dustbin-augmented transport path plus a Sinkhorn-only Pfam training harness with the reported default settings. It is intended to reproduce and inspect the core kernel/training artifact without cloud launchers or private storage paths. The appendix protocol notes specify the reported run parameters, TPU type, objective settings, and external assets used for the experiments.

\item {\bf Experimental setting/details}
    \item[] Question: Does the paper specify all the training and test details needed to understand the results?
    \item[] Answer: \answerYes{}
    \item[] Justification: The validation section and experimental protocol notes specify the masked validation settings, the four-config TPU screen setup, the promoted-run objective weights, the TPU type, and the main throughput and checkpoint milestones. The supplementary kernel and Sinkhorn-only training harness match the transport path analyzed in the paper.

\item {\bf Experiment statistical significance}
    \item[] Question: Does the paper report error bars suitably and correctly defined or other appropriate information about the statistical significance of the experiments?
    \item[] Answer: \answerNo{}
    \item[] Justification: The empirical evidence is a methods-focused exactness suite, a three-seed local bias-certificate check, an orbit-factorization check, a short TPU screen, and a single promoted run. We do not report multi-seed confidence intervals or other formal significance summaries for the TPU runs.

\item {\bf Experiments compute resources}
    \item[] Question: For each experiment, does the paper provide sufficient information on the computer resources (type of compute workers, memory, time of execution) needed to reproduce the experiments?
    \item[] Answer: \answerYes{}
    \item[] Justification: The paper identifies the TPU v6e-8 platform, reports the short-screen and promoted-run budgets, gives throughput measurements, and states the masked validation scales used for exactness checks.

\item {\bf Code of ethics}
    \item[] Question: Does the research conducted in the paper conform, in every respect, with the NeurIPS Code of Ethics?
    \item[] Answer: \answerYes{}
    \item[] Justification: The work is a methods and systems paper on differentiable attention kernels and benchmark training. It does not involve crowdsourcing, human subjects, or prohibited applications, and the manuscript does not make deployment claims beyond the reported experimental scope.

\item {\bf Broader impacts}
    \item[] Question: Does the paper discuss both potential positive societal impacts and negative societal impacts of the work performed?
    \item[] Answer: \answerYes{}
    \item[] Justification: The discussion limits the claims to factorization-backed trainability and notes the main practical considerations: compute cost, scope of the surrogate, and the difference between reconstruction-optimized and alignment-diagnostic quantities.

\item {\bf Safeguards}
    \item[] Question: Does the paper describe safeguards that have been put in place for responsible release of data or models that have a high risk for misuse (e.g., pretrained language models, image generators, or scraped datasets)?
    \item[] Answer: \answerNA{}
    \item[] Justification: The submission does not release a high-risk generative model or a new scraped dataset; it studies transport-based attention kernels and benchmark training methodology.

\item {\bf Licenses for existing assets}
    \item[] Question: Are the creators or original owners of assets used in the paper properly credited and are the license and terms of use explicitly mentioned and properly respected?
    \item[] Answer: \answerYes{}
    \item[] Justification: The protocol notes identify the public Pfam data and the public ESM-2 weights used in the experiments and state the license information included in the submission text: CC0 for Pfam/InterPro and MIT for the ESM-2 checkpoint used for embeddings.

\item {\bf New assets}
    \item[] Question: Are new assets introduced in the paper well documented and is the documentation provided alongside the assets?
    \item[] Answer: \answerNA{}
    \item[] Justification: The submission does not introduce a new public dataset or model asset release.

\item {\bf Crowdsourcing and research with human subjects}
    \item[] Question: For crowdsourcing experiments and research with human subjects, does the paper include the full text of instructions given to participants and screenshots, if applicable, as well as details about compensation (if any)?
    \item[] Answer: \answerNA{}
    \item[] Justification: The paper does not involve crowdsourcing or human-subject experiments.

\item {\bf Institutional review board (IRB) approvals or equivalent for research with human subjects}
    \item[] Question: Does the paper describe potential risks incurred by study participants, whether such risks were disclosed to the subjects, and whether Institutional Review Board (IRB) approvals (or an equivalent approval/review based on the requirements of your country or institution) were obtained?
    \item[] Answer: \answerNA{}
    \item[] Justification: The paper does not involve human-subject research.

\item {\bf Declaration of LLM usage}
    \item[] Question: Does the paper describe the usage of LLMs if it is an important, original, or non-standard component of the core methods in this research? Note that if the LLM is used only for writing, editing, or formatting purposes and does \emph{not} impact the core methodology, scientific rigor, or originality of the research, declaration is not required.
    \item[] Answer: \answerNA{}
    \item[] Justification: LLMs are not part of the core method, proofs, experimental protocol, or scientific claims in this submission.

\end{enumerate}
